\scanpage{64}
\mspage{13}

conditions $\textcircled{1}$ à $\textcircled{5}$ précédentes.

% slanted
\margin{cf axiome $W(6)$ page suivante}

Bien sûr, on aimerait un axiome 6 plus satisfaisant, qui joint aux précédents assure que $W=W_\infty$.$^*$ Ou est-ce que ces propriétés assurent déjà que $W=W_\infty$ ? Pour le voir, il faudrait peut-être que j'en sache plus sur la structure des objets $W_\infty$-asphériques ? Si c'était faux, je songe à un axiome de ``descente homotopique'', qui indiquerait qu'un objet $W_\infty$-asphérique est aussi $W$-asphérique.

Je note dès maintenant que si des $W_i$, $i\in I$, satisfont aux axiomes $1^\circ$ à $5^\circ$, il en est de $=$ des $\bigcap W_i$. Donc il y a un plus petit ens. de flèches $W_{\unsure{\omega|\infty}}$ satisfaisant à ces conditions. S'il \unsure{m'est|m'était|était} égal à $W_\infty$, i.e. si les axiomes 1 à 5 ne caractérisent déjà $W_\infty$, ce serait un \deleted{\unsure{résultat}} localiseur vraiment très intéressant !

% slanted
\margin{C'est vrai pour tout ens. d'axiomes dans \unsure{la catégorie|le catalogue} d'axiomes de stabilité}

Je me rends compte que je n'ai utilisé qu'une forme assez faible de l'axiome $\textcircled{4}$ -- essentiellement, la validité de la prop. 4. Je vais résumer différentes variantes de l'axiome de fibration.

$\textcircled{\underline{4}}$ a$_1$) Si $f$\unsure{$:$|$\in$}$X\to Y$ dans $\mathrm{Cat}$ est une $W$-fibration gauche, c'est une $W$-fibration \inserted{ultra} forte

a$_2$) Si $f:X\to Y$ est une $W$-fibration droite, c'est une $W$-fibration \inserted{ultra} forte

\deleted{\unsure{b|d}}a$'$) Si $f:X\to Y$ est \inserted{$W$-parfait, i.e.} à la fois $W$-propre et $W$-lisse$^*$, alors $f$ est une $W$-fibration \inserted{ultra} forte

% slanted
\margin{ce qui implique que c'est une $W$-fibration gauche \uline{et} droite}

a$''$) Si $f:X\to Y$ est à la fois homotopiquement propre et homotopiquement lisse (cf NB page 10),$^{\dagger}$ c'est une $W$-fibration \inserted{ultra} forte. $^{\dagger}$ Il vaut mieux dire : si $f$ est $W_\infty$-parfait, c'est une $W$-fibr. ultra-forte

\deleted{\unsure{d|c}}b) Soit $f:\unsure{\alpha|X}\to Y$ une flèche de $\mathrm{Cat}_{/S}$, avec $X$ et $Y$ des $W$-fibrés sur $S$. Si $f\in W$, alors $\forall s\in S$, $f_s:X_s\unsure{\mathrel{-}|\to}Y_s$ est $\in W$.

\deleted{\unsure{d|c}}b$'$) Soit $f:X\to Y$ une flèche de $\mathrm{Cat}_{/S}$, qui soit \deleted{à la fois} $W$-parfaite et $W_\infty$-parfaite, \deleted{\unsure{et si}} alors $f\in\unsure{V|W}$ implique que les fibres de $f$ sont $W$-asphériques. (Il vaut mieux dire : si $f$ est $W_\infty$-parfait, il est une $W$-fib. ultra-forte)

% archive sheet number: 64
